\section{Verification with \GReusot}
\label{sec:grust:greusot}

This section presents \GReusot, a \GRust wrapper around the \Creusot deductive verifier. It begins
with an overview of safety considerations in the automotive industry, then introduces \GReusot's
syntax, and finally illustrates its usage and limitations through an example.

\subsection{Safety Properties in Automotive}
\label{sec:grust:greusot:auto}

The \iso standard defines the functional safety requirements for automotive electrical/electronic
systems. It introduces Automotive Safety Integrity Levels (\Asil), ranging from \Asil~A (least
critical, \eg blinking indicators) to \Asil~D (most critical, \eg power steering). Hazards without
direct safety impact are classified as Quality Management (QM). Although QM systems require no
formal safety measures, they strongly influence the user experience.

Critical systems must avoid common faults such as uninitialized memory access, invalid pointer
dereferences, arithmetic overflow, division by zero, and memory leaks. Because runtime checks are
often costly or impractical in embedded environments, static verification is particularly important
for \Asil~D systems.

\GRust compiles to \Rust, a language that combines strong memory safety with low-level efficiency.
\Rust enforces compile-time restrictions on shared mutation, thereby preventing undefined behavior.
However, some patterns---such as mutexes---require shared mutability and cannot be expressed
entirely in \emph{safe} \Rust. To support them, \Rust provides encapsulated \emph{unsafe} code
through restricted APIs. The \RustBelt project~\cite{DBLP:journals/pacmpl/0002JKD18}
establishes a framework to prove the absence of undefined behavior for such APIs, including a
verified mutex library.

\GRust offers two kernels: a synchronous one (\GRSync) and an event-driven one (\GRFRP). The former
is fully compatible with safe \Rust, while the latter abstracts bus operations \via the
\texttt{futures} library~\cite{futures}, which currently lacks verification in \RustBelt.

The deductive verifier \Creusot~\cite{DBLP:conf/icfem/DenisJM22} builds on \Rust's memory safety to
prove functional properties with relatively few annotations (compared to those required for \Ci
code). It translates annotated safe \Rust programs into
\textsc{Why3}~\cite{DBLP:conf/esop/FilliatreP13}, which in turn invokes automated solvers such as
Z3~\cite{DBLP:conf/tacas/MouraB08}, CVC4~\cite{DBLP:conf/cav/BarrettCDHJKRT11}, or
Alt-Ergo~\cite{conchon2018alt}. By default, \Creusot checks for the absence of runtime panics,
including arithmetic overflows, across the program. Developers can also specify custom properties.
To integrate these capabilities, \GRust provides \GReusot, a wrapper of \Creusot's specification
language, for verifying \GRSync components and functions.

\subsection{Syntax}
\label{sec:grust:greusot:syntax}

\Cref{fig:syntax:greusot} shows the EBNF grammar of \GReusot, which extends \GRSync declarations
with formal specifications embedded in their signatures.

\begin{figure}[!ht]
  \centering
  \begin{bnf}(comment = {:is:})
    \FunItem ::= \funItem{\ID}{\many{(\ID \: \kft{:} \: \TY)}}{\TY}{\many{\Decl}} ;;
    \CompItem ::= \compItem{\ID}{\many{(\ID \: \kft{:} \: \STY)}}{\many{(\ID \: \kft{:} \: \STY)}}{\many{\Eq}} ;;
    \Spec ::= \many{\Clause} ;;
    \Clause ::= \reqClause{\Term} // \ensClause{\Term} ;;
    \Term ::= \CST // \ID // \lastExpr{\ID} // \opExpr{\many{\Term}} // \tuple{\Term} // \appExpr{\ID}{\many{\Term}}
    | \implTerm{\Term}{\Term} // \forallTerm{\ID}{\TY}{\Term}
    | \whenTerm{\EPat}{\Term} // \resTerm ;;
  \end{bnf}
  \caption[EBNF grammar of \GReusot specifications.]
  {EBNF grammar of \GReusot specifications in component and function declarations.
    Terminal symbols are in purple.
    Repetitions of X are denoted \many{\text{X}}.
    \ID are identifiers composed of ASCII characters.
    \CST are constants.
    \OP represents $n$-ary operators (\eg conjunction or conditional).}
  \label{fig:syntax:greusot}
\end{figure}

A specification (\Spec) consists of clauses (\many{\Clause}) defining preconditions
(\kft{requires}) and postconditions (\kft{ensures}). A \kft{requires} clause states the assumptions
under which the \kft{ensures} clause must always hold. Clause properties are expressed as terms,
which may be constants (\CST), identifiers (\ID), memories of past executions (\lastExpr{\ID}),
$n$-ary operations (\opExpr{\many{\Term}}), tuples (\tuple{\Term}), function calls
(\appExpr{\ID}{\many{\Term}}), implications (\implTerm{\Term}{\Term}), universal quantifiers
(\forallTerm{\ID}{\TY}{\Term}), event-based implications (\whenTerm{\EPat}{\Term}), or the return
value of the specified function (\resTerm).

\subsection{Compiling \GReusot Specifications}
\label{sec:grust:greusot:compil}

Compilation of \GRust transforms \kft{function} and \kft{component} declarations into \Rust
functions, encoding their respective functions or state machines. Logical variants of functions are
also generated for use in specifications. \GReusot contracts are transformed into equivalent
\Creusot clauses in the output \Rust code, expressed through \Creusot's \texttt{\#[requires(...)]}
and \texttt{\#[ensures(...)]} macros.

Verification is then performed by running \texttt{cargo creusot prove} on the generated code. Since
the \GRust compiler is implemented as a \Rust procedural macro, \Creusot can be applied directly to
the compiled crate of a \GRust program. Verification conditions are discharged by \Creusot through
several SMT solvers.

\subsection{Specifying the ACC in \GReusot}
\label{sec:grust:greusot:example}

The \verb|safety_distance| function and the \verb|acc| component from the ACC example
(\Cref{sec:grust:example}) are extended with \GReusot specifications in
\Cref{lst:grust:greusot:fun:spec,lst:grust:greusot:comp:spec}. Since \Creusot does not support
floating-point numbers, values are expressed as integers.

The function \verb|safety_distance| is intended for subject vehicle speeds between 0 and 50 meters
per second (Line 2). The front vehicle must be slower, with a maximum speed difference of 10 meters
per second (Line 3). Under these assumptions, Line 4 guarantees that the computed safety distance is
always positive, does not exceed 150 meters (the radar detection range), and cannot overflow.

\begin{lstlisting}[
      label=lst:grust:greusot:fun:spec,
      language=GRust, 
      caption={\GReusot specification of the \texttt{safety\_distance} function from 
      \Cref{fig:code_acc}.}, captionpos=b
]
function safety_distance(sv_v: int, fv_v: int) -> int
  requires { 0 < sv_v && sv_v <= 50 }
  requires { 0 < fv_v && fv_v < sv_v && sv_v - fv_v <= 10 }
  ensures  { 0 < result && result < 150 }
{
  let (rho: int, b_max: int) = (1, 6);
  let sv_d_stop: int = sv_v*rho + sv_v^2/(2*b_max);
  let fv_d_stop: int = fv_v^2/(2*b_max);
  return sv_d_stop - fv_d_stop;
}
\end{lstlisting}

The generated code corresponding to this specification appears in
\Cref{lst:grust:greusot:fun:codegen}. Because \GRSync functions are nearly identical to their
\Rust counterparts (\Cref{sec:grust:implem:choices:sync}), their specifications also translate
directly. The key point is that inside \Creusot's \texttt{requires} and \texttt{ensures} macros, we
refer to the \emph{logical representation} of types (\eg \texttt{Int} for integers instead of the
computational type \texttt{i64}, as in Line 5). In \GRust, all integers are written simply as
\texttt{int}, but after compilation the inputs and outputs must be explicitly converted to their
logical form with the \texttt{@} operator (\eg Line 1 converts \texttt{sv\_v: i64} into
\texttt{sv\_v@: Int}).

To remain compositional, \Creusot abstracts specified functions by their contracts: when a specified
function is called in another specification, only its specification (not its body) is used. To make
such calls provable, \GReusot generates a \emph{logical module} containing logical counterparts of
function bodies. These use logical types such as \texttt{Int} and are annotated with the
\texttt{\#[open]} and \texttt{\#[logic]} macros. The correspondence between a computational function
and its logical version is enforced by an \texttt{ensures} clause (Line 4).

\begin{lstlisting}[
      label=lst:grust:greusot:fun:codegen,
      language=Rust, 
      caption={Generated code for the \GReusot specification of the \texttt{safety\_distance}
      function from \Cref{lst:grust:greusot:fun:spec}.}, captionpos=b
]
#[requires(0 < sv_v@ && sv_v@ <= 50)]
#[requires(0 < fv_v@ && fv_v@ < sv_v@ && sv_v@ - fv_v@ <= 10)]
#[ensures(0 < result@ && result@ < 150)]
#[ensures(result@ == logical::safety_distance(sv_v@, fv_v@))]
pub fn safety_distance(sv_v: i64, fv_v: i64) -> i64 {
  let (rho, b_max) = (1i64, 6i64);
  let sv_d_stop = sv_v*rho + ((sv_v*sv_v) / (2*b_max));
  let fv_d_stop = (fv_v*fv_v) / (2*b_max);
  sv_d_stop - fv_d_stop
}
mod logical {
  #[open] #[logic]
  pub fn safety_distance(sv_v: Int, fv_v: Int) -> Int {
    let (rho, b_max) = (1, 6);
    let sv_d_stop = sv_v*rho + ((sv_v*sv_v) / (2*b_max));
    let fv_d_stop = (fv_v*fv_v) / (2*b_max);
    sv_d_stop - fv_d_stop
  }
}
\end{lstlisting}

The specification of the \verb|acc| component
(\Cref{lst:grust:greusot:comp:spec}) assumes at Lines 5-7 that, whenever the activation condition
\verb|c| holds:
\begin{itemize}
  \item the subject vehicle's speed lies in the valid range of \verb|safety_distance|,
  \item the derivative of the inter-vehicle distance is negative (vehicles are approaching) and
        remains within the bounds of \verb|safety_distance|,
  \item the distance between the vehicles is above the minimum threshold for maximum braking.
\end{itemize}
Additionally, it assumes that the radar distance \verb|d| does not exceed 150 meters (Line 3). The
postconditions guarantee absence of division by zero and overflow, and that the braking rate always
remains within bounds (Line 9). Proving the specification of \verb|acc| requires an auxiliary
function \verb|compute_braking| to guide the solvers (Lines 24-33).

\begin{lstlisting}[
  label=lst:grust:greusot:comp:spec,
  language=GRust, 
  caption={\GReusot specification of the \texttt{acc} component from \Cref{fig:code_acc}.},
  captionpos=b
]
component acc(c: bool, d: int, v: int, s: int) -> (b: int)
  // radar detection limitation
  requires { d < 150 }
  // ACC scope of usage
  requires { c => (0 < s && s <= 50) && (0 < s+v && v < 0 && -v <= 10) }
  // there is enough distance to brake at maximum rate
  requires { c => d - safety_distance(s, s+v) > (v^2)/12 }
  // braking rate is in correct interval
  ensures  { 0 <= b && b <= 6 }
{
  match c {
    true => {
      b = compute_braking(d - d_safe, v);
      let d_safe: int = safety_distance(s, fv_v);
      let fv_v: int = s + v;
    }
    false => {
      b = 0;
      let (fv_v: int, d_safe: int) = (0, 0);
    }
  }
}
// Intermediate braking function to pass the proof
function compute_braking(d_grace: int, v: int) -> int
  // scope
  requires { (0 < d_grace &&  d_grace < 150) && (v < 0 && -v <= 10) }
  // there is enough distance to brake at maximum rate
  requires { d_grace > (v^2)/12 }
  // braking rate is in correct interval
  ensures  { 0 <= result && result <= 6 }
{
  return (v^2) / (2*d_grace);
}
\end{lstlisting}

The generated code is shown in \Cref{lst:grust:greusot:comp:codegen}. As explained in
\Cref{sec:grust:implem:choices:sync}, the prototype \GRust compiler produces for each component an
implementation of the \texttt{Component} trait from the \GRust core library, including
\texttt{init} and \texttt{step} functions. The \texttt{init} function builds the initial state,
while the \texttt{step} function takes the current inputs and a mutable reference to the state,
computes outputs, and updates the state. An access to a memory \lastExpr{x} is compiled into a field
access on \texttt{state}, referring to its value \emph{before} the computations.

The call to \verb|safety_distance| in the \verb|acc| specification is compiled to a call to its
logical counterpart \verb|logical::safety_distance| (Line 4), which preserves the relationship
between inputs (\verb|input.s@| and \verb|input.s@ + input.v@|) and the result. Similarly, the
intermediate function \verb|compute_braking| generates a logical version in the \verb|logical|
module (Lines 30--35).

\begin{lstlisting}[
  label=lst:grust:greusot:comp:codegen,
  language=Rust, 
  caption={Generated code for the \GReusot specification of the \texttt{acc} component from
  \Cref{lst:grust:greusot:comp:spec}.}, captionpos=b
]
#[requires(input.d@ < 150)]
#[requires(input.c ==> (0 < input.s@ && input.s@ <= 50) && 
                        (0 < input.s@ + input.v@ && input.v@ < 0 && - input.v@ <= 10))]
#[requires(input.c ==> input.d@ - logical::safety_distance(input.s@, input.s@ + input.v@)
                        > (input.v@*input.v@) / 12)]
#[ensures(0 <= result@ && result@ <= 6)]
fn step(state: &mut AccState, input: AccInput) -> i64 {
  let (d_safe, b, fv_v) = match input.c {
    true => {
      let fv_v = input.s + input.v;
      let d_safe = safety_distance(input.s, fv_v);
      let b = compute_braking(input.d - d_safe, input.v);
      (d_safe, b, fv_v)
    }
    false => {
      let b = 0i64;
      let (fv_v, d_safe) = (0i64, 0i64);
      (d_safe, b, fv_v)
    }
  };
  b
}
#[requires((0 < d_grace@ && d_grace@ < 150) && (v@ < 0 && - v@ <= 10))]
#[requires(d_grace@ > (v@*v@) / 12)]
#[ensures(0 <= result@ && result@ <= 6)]
#[ensures(result@ == logical::compute_braking(d_grace@, v@))]
pub fn compute_braking(d_grace: i64, v: i64) -> i64 {
  (v*v) / (2i64*d_grace)
}
mod logical {
  #[open] #[logic]
  pub fn compute_braking(d_grace: Int, v: Int) -> Int {
    (v*v) / (2*d_grace)
  }
}
\end{lstlisting}

\subsection{Future Work on \GReusot}
\label{sec:grust:greusot:future}

Formal verification is mandatory for certain \Asil levels. We developed \GReusot as a proof of
concept to demonstrate the feasibility of formal verification in \GRust, thereby strengthening its
relevance in the automotive domain. However, since the development of such a tool was not our
initial objective, several improvements remain for future work. This section outlines a
non-exhaustive list of enhancements that would significantly improve the usability of \GReusot.

Deductive verification often requires the programmer to guide the solvers, as illustrated in our
example with the auxiliary function \verb|compute_braking| (\Cref{lst:grust:greusot:comp:spec}).
Programmers could be better supported by enabling them to insert assertions on local variables. In
functions, such assertions could be placed directly before the declaration of the variable. In
components, since equations are unordered, assertions could be added at the beginning of the set of
equations. This avoids polluting the signature with irrelevant details, as assertions concern only
local variables. In the generated \texttt{step} function, these assertions could then be compiled
into \Creusot assertions inserted after the creation of the relevant identifiers.

Another limitation concerns the use of memories (\lastExpr{\ID}) in component specifications, which
are currently not compositional. To verify a component that calls another, the called component must
be provided with a specification. If this specification refers to internal memory in a
\kft{requires} clause, the caller must establish the condition. However, the caller cannot access
the internal memory of another component. This usability issue could be addressed by restricting
properties on memories to \kft{invariant} clauses. These invariants would be checked as
\kft{ensures} clauses in the \texttt{init} function, as \kft{requires} clauses referring on the
\texttt{state} value \emph{before} the computations, and as \kft{ensures} clauses referring on the
\texttt{state} value \emph{after} the computations (using \verb|^state| instead of \texttt{state}).
To make invariants usable in practice, we would also need to permit the use of inputs in the
\texttt{init} function.

A further limitation is the treatment of component calls within specifications. Each component call
has its own state, and these states do not interfere with one another. If a specification involves a
component call, two options arise: either we create a ghost state for the component, without any
\apriori knowledge of its contents; or we reuse an existing state and call a logical version of its
\texttt{step} function (analogous to function calls), which does not modify the state. In the latter
case, care must be taken to select the component state that corresponds to a call with the same
inputs as in the specification. The model checker Kind 2~\cite{DBLP:conf/cav/ChampionMST16} supports
specifications for Lustre nodes using the k-induction
technique~\cite{DBLP:conf/fmcad/SheeranSS00}, a well-established method for verifying transition
systems such as stateful nodes (or components). Integrating a similar k-induction approach into
\Creusot has not yet been investigated.
